\documentclass[11pt]{article}

\usepackage[a4paper,margin=1in]{geometry}
\usepackage{amsmath}
\usepackage{amssymb}
\usepackage{bm}
\usepackage{booktabs}
\usepackage{array}
\usepackage{graphicx}
\usepackage{enumitem}

\title{Initial Overview for Implementation:\\
Temporal Variational Autoencoder and Attention-Based Transformer for Sequential Equity Return Prediction}

\author{Aarush Ram Anandh}
\date{\today}

\begin{document}

\maketitle

\tableofcontents

\newpage

\section{Problem Definition}

We are given panel time-series financial data consisting of multiple companies observed across multiple months.

Each observation contains:
\begin{itemize}
    \item company identifier
    \item timestamp
    \item accounting variables
    \item market variables
    \item engineered financial factors
    \item lagged historical statistics
\end{itemize}

The objective is to predict the next-month gross return for a company using its complete historical trajectory up to the current month.

Formally, for company $i$:

\[
(X_{i,1}, X_{i,2}, \dots, X_{i,t})
\rightarrow
r_{i,t+1}
\]

where:
\begin{itemize}
    \item $X_{i,t}$ denotes the feature vector at month $t$
    \item $r_{i,t+1}$ denotes the next-month gross return
\end{itemize}

The model therefore learns a sequential mapping rather than an independent row-wise regression function.

\section{Dataset Structure}

Each row in the dataset corresponds to:

\[
(\text{company}, \text{month}, \text{features})
\]

Example feature groups include:
\begin{itemize}
    \item market capitalization
    \item price
    \item book-to-market ratio
    \item profitability
    \item investment factor
    \item lagged returns
    \item momentum
    \item accounting statistics
\end{itemize}

The dataset is therefore:
\begin{itemize}
    \item multivariate
    \item temporal
    \item panel structured
\end{itemize}

\section{Data Preprocessing}

\subsection{Chronological Sorting}

For every company:

\[
D_i =
\{
(X_{i,1}, r_{i,1}),
(X_{i,2}, r_{i,2}),
\dots,
(X_{i,T_i}, r_{i,T_i})
\}
\]

the observations must be sorted chronologically.

\subsection{Target Construction}

The prediction target is the next-month return.

For every timestep:

\[
y_{i,t} = r_{i,t+1}
\]

Thus:
\begin{itemize}
    \item input uses information available at time $t$
    \item target corresponds to month $t+1$
\end{itemize}

\subsection{Sequence Construction}

A fixed rolling window of length $L$ is used.

For example, if:
\[
L = 12
\]

then each training sample becomes:

\[
S_{i,t}
=
(X_{i,t-11},
X_{i,t-10},
\dots,
X_{i,t})
\]

with target:

\[
y_{i,t} = r_{i,t+1}
\]

Thus:
\[
S_{i,t} \in \mathbb{R}^{L \times d}
\]

where:
\begin{itemize}
    \item $L$ is sequence length
    \item $d$ is feature dimension
\end{itemize}

\subsection{Sliding Window Construction}

The rolling window slides across time.

Example:

\[
(X_1,\dots,X_{12}) \rightarrow r_{13}
\]

\[
(X_2,\dots,X_{13}) \rightarrow r_{14}
\]

\[
(X_3,\dots,X_{14}) \rightarrow r_{15}
\]

This converts the panel dataset into a supervised sequential learning dataset.

\section{Train / Validation / Test Split}

Random splitting must never be used because it causes temporal leakage.

Instead, chronological splitting is performed.

Example:

\begin{center}
\begin{tabular}{ll}
\toprule
Split & Time Range \\
\midrule
Training & 2000 -- 2018 \\
Validation & 2019 -- 2021 \\
Testing & 2022 -- 2025 \\
\bottomrule
\end{tabular}
\end{center}

This ensures:
\begin{itemize}
    \item the model never sees future data during training
    \item evaluation remains realistic
    \item deployment conditions are preserved
\end{itemize}

\section{Feature Normalization}

Financial variables have different scales.

Examples:
\begin{itemize}
    \item market capitalization may be extremely large
    \item momentum may be centered near zero
    \item profitability ratios may contain small decimals
\end{itemize}

Normalization is therefore required.

For each feature:

\[
x' =
\frac{x - \mu}{\sigma}
\]

where:
\begin{itemize}
    \item $\mu$ is training mean
    \item $\sigma$ is training standard deviation
\end{itemize}

Important:
\begin{itemize}
    \item normalization statistics must be computed only on training data
    \item validation and test sets must reuse training statistics
\end{itemize}

\section{Categorical Features}

Categorical labels such as:
\begin{itemize}
    \item momentum labels
    \item size labels
    \item book-to-market labels
\end{itemize}

must be encoded numerically.

Possible methods:
\begin{itemize}
    \item one-hot encoding
    \item learned embeddings
\end{itemize}

Embedding layers are preferred for transformer architectures.

\section{Temporal Variational Autoencoder}

\subsection{Objective}

The Temporal VAE learns compressed latent representations of historical financial sequences.

Instead of directly predicting returns, the model first learns a latent state:

\[
z_t
=
f(X_{1:t})
\]

The latent vector should capture:
\begin{itemize}
    \item market regimes
    \item momentum structure
    \item volatility patterns
    \item long-term financial behavior
\end{itemize}

\subsection{Encoder}

The encoder receives the sequence:

\[
S_{i,t}
=
(X_{i,t-L+1}, \dots, X_{i,t})
\]

and produces:

\[
q_\phi(z|S)
=
\mathcal{N}(\mu, \sigma^2)
\]

The encoder outputs:
\begin{itemize}
    \item latent mean $\mu$
    \item latent variance $\sigma^2$
\end{itemize}

\subsection{Latent Sampling}

Using the reparameterization trick:

\[
z = \mu + \sigma \odot \epsilon
\]

where:

\[
\epsilon \sim \mathcal{N}(0, I)
\]

\subsection{Decoder}

The decoder reconstructs the original sequence:

\[
\hat S = g_\theta(z)
\]

The objective is to preserve sequential structure inside the latent space.

\subsection{Loss Function}

The VAE minimizes:

\[
\mathcal{L}
=
\mathcal{L}_{recon}
+
\beta D_{KL}(q(z|x)\|p(z))
\]

where:
\begin{itemize}
    \item reconstruction loss preserves information
    \item KL divergence regularizes latent space
\end{itemize}

Typical reconstruction loss:

\[
\mathcal{L}_{recon}
=
\|S - \hat S\|^2
\]

\section{Transformer Forecasting Model}

\subsection{Objective}

The transformer predicts next-month return from temporal latent representations.

Input:

\[
(z_1, z_2, \dots, z_t)
\]

Output:

\[
\hat r_{t+1}
\]

\subsection{Positional Encoding}

Transformers do not naturally understand temporal order.

Therefore positional encodings are added:

\[
h_t = z_t + p_t
\]

where:
\begin{itemize}
    \item $z_t$ is latent representation
    \item $p_t$ is temporal positional encoding
\end{itemize}

\subsection{Self-Attention}

Attention computes relationships between months.

For query $Q$, key $K$, and value $V$:

\[
\text{Attention}(Q,K,V)
=
\text{softmax}
\left(
\frac{QK^T}{\sqrt{d_k}}
\right)V
\]

This allows the model to learn:
\begin{itemize}
    \item long-range dependencies
    \item regime persistence
    \item historical relevance
\end{itemize}

\subsection{Prediction Head}

The final transformer representation:

\[
h_t
\]

is passed into a regression head:

\[
\hat r_{t+1}
=
W h_t + b
\]

\section{Training Pipeline}

\subsection{Step 1: Sequence Generation}

Construct rolling windows for every company.

Output tensor:

\[
X \in \mathbb{R}^{N \times L \times d}
\]

where:
\begin{itemize}
    \item $N$ = number of training samples
    \item $L$ = sequence length
    \item $d$ = feature dimension
\end{itemize}

\subsection{Step 2: Train Temporal VAE}

The VAE is trained first.

Objective:
\begin{itemize}
    \item compress sequences
    \item learn smooth latent manifold
\end{itemize}

Output:
\[
z_t
\]

for every sequence.

\subsection{Step 3: Transformer Training}

Latent sequences are fed into the transformer.

Objective:
\[
\hat r_{t+1}
\approx
r_{t+1}
\]

Typical regression loss:

\[
\mathcal{L}_{pred}
=
(r_{t+1} - \hat r_{t+1})^2
\]

\section{Inference Procedure}

At inference time:

\begin{enumerate}
    \item collect most recent $L$ months
    \item normalize using training statistics
    \item encode sequence using VAE encoder
    \item obtain latent sequence
    \item feed latent sequence into transformer
    \item predict next-month return
\end{enumerate}

Mathematically:

\[
(X_{t-L+1}, \dots, X_t)
\rightarrow
(z_{t-L+1}, \dots, z_t)
\rightarrow
\hat r_{t+1}
\]

\section{Evaluation Metrics}

Possible evaluation metrics include:

\begin{itemize}
    \item Mean Squared Error
    \item Mean Absolute Error
    \item Directional Accuracy
    \item Rank Correlation
    \item Sharpe Ratio
\end{itemize}

Financial prediction quality is often more dependent on ranking quality than raw numerical precision.

\section{Expected Research Contributions}

The project investigates whether:
\begin{itemize}
    \item latent temporal representations improve forecasting
    \item transformers capture long-range financial dependencies
    \item attention mechanisms improve interpretability
    \item probabilistic latent modeling stabilizes prediction
\end{itemize}

\section{Future Extensions}

Possible future directions:
\begin{itemize}
    \item regime-aware transformers
    \item graph neural networks across companies
    \item reinforcement learning portfolio allocation
    \item diffusion-based financial sequence generation
    \item contrastive representation learning
\end{itemize}

\end{document}
